\subsection{Panel de control (Dashboard) y gestión de la convivencia}
\label{subsec:dashboard_convivencia}

El panel de control principal representa el centro neurálgico de la experiencia diaria del estudiante o inquilino en Colivi. Tras autenticarse, la plataforma presenta una vista consolidada que centraliza el estado actual de la convivencia residencial, los compromisos económicos compartidos y la distribución de responsabilidades domésticas.

\subsubsection{Estructura de la interfaz y composición de widgets}

En el frontend, el panel de control se implementa a través del componente de página principal (\texttt{HomePage.tsx}) y componentes de agregación modular en \texttt{components/layout/}:

\begin{itemize}
    \item \textbf{Encabezado de comunidad y membresía activa}: Muestra el nombre del hogar actual (por ejemplo, \emph{«Piso Estudiantes Centro»}), la insignia identificativa de rol del usuario (\texttt{ADMIN} o \texttt{MEMBER}) y la lista interactiva de compañeros de piso con sus respectivos avatares y códigos de color personalizados, facilitando la distinción visual instantánea.
    \item \textbf{Widget de balance financiero rápido}: Sintetiza el saldo neto del usuario en el piso. Si el balance es positivo, indica el importe total que sus compañeros le adeudan; si es negativo, alerta en tono de énfasis sobre los pagos pendientes, ofreciendo un botón de acceso directo a la liquidación.
    \item \textbf{Widget de tareas domésticas inminentes}: Expone las labores asignadas al usuario para la semana en curso (como la limpieza de zonas comunes o retirada de residuos), permitiendo marcar tareas como finalizadas con un solo clic o solicitar un «rescate» gamificado a la comunidad.
    \item \textbf{Gestión de invitaciones del hogar}: Para los administradores de la vivienda, el panel proporciona un botón para copiar al portapapeles o regenerar el código de invitación alfanumérico (ejemplo: \texttt{PISO-CENTRO-2024}), permitiendo la adhesión inmediata de nuevos inquilinos sin requerir trámites administrativos complejos.
\end{itemize}

\subsubsection{Agregación de datos y optimización de consultas en backend}

Para evitar el problema clásico de sobrecarga por múltiples peticiones HTTP en cascada (*Waterfall Requests*), el backend expone un endpoint consolidado en \texttt{HomeController} (\texttt{GET /api/v1/homes/{homeId}/summary}):

\begin{enumerate}
    \item \textbf{Consolidación en un único DTO (\texttt{HomeSummaryDTO})}: El servicio \texttt{HomeService} orquesta en una única transacción de sólo lectura la recopilación de los miembros activos del hogar, los últimos tres gastos registrados, el balance contable individual del solicitante y el cuadrante de tareas activas.
    \item \textbf{Prevención de consultas N+1 en Hibernate}: Se diseñaron consultas especializadas en \texttt{HomeRepository} empleando cláusulas \texttt{LEFT JOIN FETCH} sobre las entidades relacionadas (\texttt{members}, \texttt{expenses}, \texttt{chores}), asegurando que la base de datos PostgreSQL resuelva toda la información en un número acotado de consultas SQL altamente optimizadas.
    \item \textbf{Sincronización reactiva con TanStack Query}: En el cliente, la consulta del panel se suscribe a una clave de caché unificada (\texttt{['home', homeId, 'summary']}). Cuando cualquier miembro del piso añade un gasto o completa una tarea, las mutaciones invalidan selectivamente esta clave, provocando un refresco en segundo plano que mantiene la interfaz siempre actualizada sin saltos visuales ni parpadeos.
\end{enumerate}

